%قبل از شروع مسئله ترجیحاً چیزی ننویسید. اما اگر خواستید بنویسید، در قالب لیستی (فقرات و شمارش) بنویسید.
%برای اضافه کردن نمره، بعدی از شروع‌مسئله، [x] را اضافه کنید و به جای x  نمره مورد نظر را قرار دهید. برای زیرمسئله‌ها هم به همین صورت است.
%\شروع‌مسئله[10]
\شروع‌مسئله [20]
%\متن‌سیاه{سوال عملی اول:}\\
%\متن‌سیاه{دراپ‌تست}\\
درختی \متن‌لاتین{n} راسی و ریشه دار داریم. رئوس با اعداد ۱ تا \متن‌لاتین{n} شماره گذاری شده‌اند و راس شماره ۱ ریشه‌ی درخت است. علاوه بر عدد، هر راس یک مقدار ویژه هم دارد. به ازای یک راس تمام مقدار ویژه های رئوس زیر درختش را در نظر بگیرید. مینیمم تفاوت بین مقدار ویژه‌ها، امتیاز ویژه‌ی آن راس است. حال روشی ارائه دهید که بتوانیم امتیاز ویژه‌ی همه‌ی راس‌ها را در 
$O(n\log(n)^2)$
بدست آوریم.

\solution{
پاسخ. برای این کار، تابعی به شکل 
$f(v)$
تعریف می‌کنیم که این تابع، یک درخت AVL ارائه می‌دهد که در آن تمام اعداد رئوس زیردرخت راس v در آن آمده است. هم چنین امتیاز ویژه‌ی راس v را نیز می‌دهد. حال می‌خواهیم روشی برای این تابع ارائه دهیم. 
برای ساخت درخت AVL ،‌ کافی است یک درخت AVL بسازیم که شامل عدد راس v است. سپس درخت‌های AVL بچه‌های راس v را بدست آوریم و به ترتیب آن‌ها را با درخت AVL راس v ادغام کنیم. 

برای پیدا کردن امتیاز ویژه هم می‌دانیم دو راس با اختلاف مینیمم، یا هر دو در زیردرخت یکی از بچه‌های راس v هستند که کافی است امتیاز ویژه‌ی آن بچه را داشته باشیم. در غیر اینصورت می‌دانیم هنگام ادغام کردن، باید مینیمم اختلاف بین اعداد دو درخت AVL را نیز بدست آوریم. 
پس برای محاسبه‌ی 
$f(v)$
ابتدا تابع
$f$
را برای همه‌ی بچه‌های راس v صدا می‌زنیم تا مواردی که در بالا نیاز است، بدست بیاید. حال برای محاسبه‌ی امتیاز ویژه و درخت AVL به صورت زیر عمل می‌کنیم. 

فرض کنیم درخت AVL راس i را با 
$AVL_i$
نشان دهیم. u را بچه‌ی v در نظر می‌گیریم و می‌خواهیم دو درخت
$AVL_v, AVL_u$
را ادغام کنیم. برای این کار، اگر 
$AVL_v$
راس های کمتری نسبت به 
$AVL_u$
داشت، تک تک اعضای 
$AVL_v$
را به 
$AVL_u$
اضافه می‌کنیم. در غیر اینصورت برعکس عمل می‌کنیم. در واقع درخت با راس‌های کمتر را به درخت با راس‌های بیشتر اضافه می‌کنیم. درخت حاصل را به راس v نسبت می‌دهیم. 

هر بار که یک عضو را اضافه می‌کنیم، برای این که ببینیم با اعداد درخت دیگر مینیمم اختلافی که می‌سازد چقدر است، بزرگترین عدد کوچکتر یا مساوی و کوچکترین عدد بزرگتش در AVL ای که به آن اضافه می‌شود هم بدست می‌آوریم. مینیمم اختلاف با این دو مقدار، مینیمم مقدار اختلافی است که می‌تواند با اعداد این درخت بسازد(چرا؟). در نتیجه وقتی زیردرخت v را با کل بچه‌هایش ادغام کردیم، امتیاز ویژه‌اش مینیمم امتیاز ویژه‌‌های بچه‌هایش و اعداد بدست آمده در ادغام هاست.

حال برای حل مسئله کافی است تنها یک بار 
$f(1)$
را صدا کنیم. در نتیجه‌ی این کار f تمام رئوس صدا زده می‌شود و امتیاز ویژه‌ی هر راسی نیز محاسبه می‌شود. حال این راه را تحلیل اردر می‌کنیم. یک راس دلخواه را در نظر بگیرید. می‌خواهیم بررسی کنیم این عدد چند بار به یک درخت AVL اضافه می‌شود. وقتی دو AVL را ادغام می کنیم، اعداد AVL کوچکتر به AVL بزرگتر اضافه می‌شوند. درنتیجه AVL نهایی از ادغام این دو AVL ، حداقل دو برابر AVL کوچکتر راس دارد. پس هر بار که یک عدد به یک درخت AVL اضافه شود، تعداد راس های AVL جدیدی که به آن اضافه می‌شود نسبت به AVL قبلی‌اش حداقل دو برابر شده است. چون هیچ AVL ای بیشتر از n راس ندارد و هر بار که یک درخت AVL به درخت AVL دیگری اضافه شود دیگر با آن کاری نداریم، می‌فهمیم عدد هر راسی حداکثر 
$\log(n)$
بار به یک درخت AVL اضافه می‌شود. هم چنین یافتن بزرگترین عدد کوچکتر یا مساوی و کوچکترین عدد بزرگتر نیز از
$O(\log(n))$
است. پس کل عملیات ها از اردر زیر است:
$$Time = O(n * \log(n) * \log(n)) = O(n{\log(n)}^2)$$

}
%\متن‌سیاه{ورودی:}\\
%در خط اول ورودی، ابتدا k و سپس m داده می‌شود.\\
%\متن‌سیاه{خروجی:}\\
%در تنها سطر خروجی، تعداد آیفون‌های مورد نیاز می‌آید. اگر در هیچ صورتی چنین کاری ممکن نبود ۱- چاپ می‌شود.\\
%\متن‌سیاه{محدودیت‌ها:}\\
%$1\leq k\leq 10^4$

%$m\leq10^9$
%\\
%\textbf{\متن‌سیاه{ورودی و خروجی نمونه}\\
%% تست اول
%\begin{latin}
%	\begin{tabularx}{\textwidth}{|X | X|}
%		\hline
%		\textbf{stdin} &
%		\textbf{stdout}\\
%		\hline
%		100 100&
%		1\\
%		\hline
%	
%	\end{tabularx}
%\end{latin}
%
%% تست دوم
%
%\begin{latin}
%	\begin{tabularx}{\textwidth}{|X | X|}
%		\hline
%		\textbf{stdin} &
%		\textbf{stdout}\\
%		\hline
%		340 20&
%		3\\
%		\hline
%		
%	\end{tabularx}
%\end{latin}}
% صورت اصلی مسئله را در این‌جا بنویسید.
%اگر مسئله چند قسمت داشت به این صورت عمل کنید:
%\شروع{زیرمسائل}
%\زیرمسئله[0] % مسئلۀ ۱آ
%\زیرمسئله % مسئلۀ ۱ب
%\زیرمسئله % مسئلۀ ۱پ
%\زیرمسئله % مسئلۀ ۱ت
%\پایان{زیرمسائل}



%%%%%%%%%%%%%%%%%%%%%%%%%%%%%%%%%%%%%%%%%%%%%%%%%%%%%
% برخی دستورات مفید برای تمرین‌ها. %
%%%%%%%%%%%%%%%%%%%%%%%%%%%%%%%%%%%%%%%%%%%%%%%%%%%%%

% فرمول‌های نمونه:
%$f(x) = \Theta(x)$
%$T(x, y) \leq \log(x) + 2^y + \binom{2n}{n}$
% {\tt code\_function}


% لیست موارد بدون شماره‌گذاری:
%\شروع{فقرات}
%    \فقره مورد اول 
%        \شروع{فقرات}
%            \فقره مورد اول داخلی 
%                \شروع{فقرات}
%                    \فقره اولین مورد خیلی داخلی 
%                    \فقره دومین مورد خیلی داخلی 
%                \پایان{فقرات}
%            \فقره مورد دوم داخلی 
%        \پایان{فقرات}
%    \فقره مورد دوم 
%\پایان{فقرات}

% لیست موارد شماره‌گذاری شده:
%\شروع{شمارش}
%    \فقره یک مورد 
%    \فقره یک مورد دیگر 
%    \فقره باز هم یک مورد دیگر
%\پایان{شمارش}

% برای نوشتن یک فرمول چند خطی که یک شماره بخوره. برای حالت‌های دیگر به سایت زیر مراجعه کنید:
% https://www.overleaf.com/learn/latex/Aligning_equations_with_amsmath
%\begin{align} 
%    \begin{split}
%        (x+y)^3 &= (x+y)^2(x+y) \\
%                &= (x^2+2xy+y^2)(x+y) \\
%                &= (x^3+2x^2y+xy^2) + (x^2y+2xy^2+y^3) \\
%                &= x^3+3x^2y+3xy^2+y^3
%    \end{split}                                 
%\end{align}

% نوشتن ماتریس و آرایه دوبعدی
%\begin{align}
%    A = 
%    \begin{bmatrix}
%        A_{11} & A_{21} \\
%        A_{21} & A_{22}
%    \end{bmatrix}
%\end{align}

% اضافه کردن تصویر با شرح. برای اضافه کردن تصویر بدون شرح، فقط دستور درج‌تصویر را استفاده کنید.
%\شروع{شکل}[h!]
%\تنظیم‌ازوسط
%  \درج‌تصویر[پهنا=0.5\پهنای‌متن]{img.jpg}
%\شرح{شرح تصویر}
%\برچسب{تصویر:۱}
%\پایان{شکل}
%\پایان{زیرمسائل}

%کد زیر را ببینید.
%%\lstinputlisting{source_filename.py}
%\begin{latin}
%\begin{lstlisting}
%# Import the modules
%import sys
%import random
%
%ans = True
%
%while ans:
%question = raw_input("Ask the magic 8 ball a question: (press enter to quit) ")
%
%answers = random.randint(1,8)
%
%if question == "":
%sys.exit()
%
%elif answers == 1:
%print "It is certain"
%
%elif answers == 2:
%print "Outlook good"
%
%elif answers == 3:
%print "You may rely on it"
%
%elif answers == 4:
%print "Ask again later"
%
%elif answers == 5:
%print "Concentrate and ask again"
%
%elif answers == 6:
%print "Reply hazy, try again"
%
%elif answers == 7:
%print "My reply is no"
%
%elif answers == 8:
%print "My sources say no"
%\end{lstlisting}
%\end{latin}
